\phantomsection
\addcontentsline{toc}{section}{Abstract}
\begin{abstract}
\sloppy\itshape
The rapid expansion of digital education and remote recruitment has increased the demand for reliable, scalable, and secure online assessment systems. Traditional examination methods often face limitations such as lack of flexibility, manual administration overhead, scalability challenges, and difficulties in ensuring fairness and candidate identity verification. This report addresses these challenges by designing and documenting a comprehensive enterprise-level online assessment management system that supports multiple stakeholders, including system administrators, enterprise administrators, staff members, candidates, and external payment services.

The proposed system provides a structured, role-based platform to manage the complete lifecycle of digital assessments. Key functionalities include enterprise registration and subscription management, staff and candidate management, exam creation and scheduling, candidate assignment, secure exam delivery, identity verification, automated grading, reporting, and certificate issuance. The system integrates dashboard-based analytics to support informed decision-making at both system-wide and enterprise-specific levels. A modular, use-case-driven design approach clearly defines system behavior, interactions, and responsibilities across different actors.

This project presents \textit{Veritas}, an AI-enhanced online assessment and monitoring system tailored for Ethiopian enterprises and institutions. In alignment with the Digital Ethiopia 2030 Strategy, Veritas emphasizes local data sovereignty through a multi-tenant architecture, enabling independent exam environments for multiple organizations. Key features include secure authentication via enterprise IDs or tokens, continuous identity verification using face recognition, intelligent proctoring (e.g., tab-switch detection, mouse tracking, IP monitoring,...), automated grading for both objective and subjective questions using NLP models, customizable dashboards for exam management and analytics, automated certificate issuance, and simulated subscription modules supporting local payments.

Developed using an Agile Scrum methodology, the prototype employs a microservices architecture (Go for performance-critical components, Python for AI modules), and has been tested for usability, security, and reliability. Evaluation demonstrates improved efficiency, reduced cheating, and enhanced trust, with AI-assisted grading minimizing human bias and bottlenecks. This work contributes to software engineering by providing a blueprint for contextually adaptable, secure assessment systems that foster technological independence in resource-constrained settings, while supporting e-learning, recruitment, and certification processes.

The significance of this work lies in its practical applicability to modern e-learning platforms, corporate recruitment, and professional certification systems. The proposed design offers a foundation for implementing robust, enterprise-ready online assessment solutions that address real-world challenges in digital evaluation environments.
\end{abstract}